\chapter*{Abstract}
\thispagestyle{empty}
\par The master's thesis entitled "Public transport monitoring system based on GPS technology" was developed by the student of the CRI-201M group, Vasile Drumea, Technical University of Moldova. This project consists of 3 chapters, introduction, conclusions and bibliography.
\par \textbf{Keywords:} IoT, MQTT, GPS, Transport, Communication, Device, Publish-Subscribe.
\par Public transport is an essential element for most people around the world. Regardless of age, occupation or financial status, people need a mode of transport on which they can rely at any time, ie the public transport system must be organized to meet demand. Users also need the most up-to-date and relevant sources of information to be able to make transport decisions. There can often be multiple options for choosing public transport for the same route, so if users have the necessary information at the right time, they can make an optimal decision based on the available transport options.
\par This project aims to study how public transport is organized and how the quality of services and access to information for potential passengers can be improved. The main purpose is to provide public transport providers with services through which vehicle information can be effectively transmitted to all customer entities within the system. For the time being, the focus of this paper is on vehicle location data, but the same mechanism could be extended to transmit more data in the future. Due to the fact that the passengers / users themselves, as well as the transport companies, may be the customers, it will be necessary to make architectural decisions that allow the integration of several types of customer entities and the scalability of the system.
